% Part: quantified-modal-logic
% Chapter: semantics
% Section: variable-domain-models

\documentclass[../../../include/open-logic-section]{subfiles}

\begin{document}

\olfileid{qml}{sem}{vdm}

\olsection{Variable-Domain Models}

\begin{explain}
Variable domain semantics combines the semantics of modal logic 
with the semantics of free logic. In this semantics, what objects
there are is allowed to vary from one world to another. Thus we have
\emph{one domain of objects per world}. A world's domain consists 
of the objects that exist at that world. The domains of two worlds
can overlap: typically, objects exist at multiple worlds. 
(Hawaii would still exist if the climate was slightly 
cooler.) 

To give world-variant and possibly overlapping domains, 
we give the model an \emph{super-domain} that gathers all 
objects that may exist at some world---and possibly includes 
some objects that exist at no world at all. Each world's domain is a 
subset of the model's outer domain.

Free logic has variants, including positive free logic
 and negative free logic. The variable-domain models
 below can be used for both positive and 
 negative free quantified modal logic. Model-wise, the only 
 difference is that models for the negative variant
 obey further constraints---constraints that ensure 
 that at a given world, !!{predicate}s only apply to objects
 existing at that world, and likewise for !!{function}s. 
 The only other difference between the two is in the 
 satisfaction clause for identity. 
\end{explain}

\begin{defn}[Variable-domain models]
    A \emph{variable-domain model} $\mModel M$ for the language $\Lang{L}$ 
    of quantified modal logic consists of the following elements:
    \begin{enumerate}
    \item \emph{Worlds:} a non-empty set~$\mWorlds{M}$
    \item \emph{Accessibility relation:} a binary relation~$\mAccess{M}$ on~$\mWorlds{M}$
    \item \emph{Super-Domain:} a non-empty set, $\mathcal{S}$
    \item \emph{Domains Function}: a function $\mDomain{M}$ that assigns
    to each world $w$ in $\mWorlds{M}$ a subset $\mDomain{M}(w)$ of $\mathcal{S}$, the \emph{domain of world $w$}
\item \emph{Interpretation of !!{constant}s:} for each !!{constant}~$c$ of
  $\Lang{L}$, !!a{element} of the super-domain, $\mAssign{c}{M} \in \mathcal{S}$
\item \emph{Interpretation of !!{predicate}s:} for each $n$-place
  !!{predicate}~$R$ of $\Lang{L}$ (other than $\eq$), a function 
  from worlds to $n$-place
  relations on the super-domain: $\mAssign{R}{M} \colon W \to \mathcal{S}^n$
\item \emph{Interpretation of !!{function}s:} for each $n$-place
  !!{function}~$f$ of $\Lang{L}$, an $n$-place function from worlds 
  to $n$-place functions on the super-domain: $\mAssign{f}{M} \colon W \to 
  (\mathcal{S}^n \to \mathcal{S})$

\end{enumerate}

We write $\mAssign{R}{M}[w]$ for the value of $\mAssign{R}{M}$ at $w$ (a $n$-place relation) and $\mAssign{f}{M}[w]$ for the value of the 
$\mAssign{f}{M}[w]$ at $w$ (a $n$-place function). 
\end{defn}

\begin{defn}[Negative variable-domain models]
  A variable-domain model is \emph{negative} if it additionally 
  obeys the following constraints:
  \begin{enumerate}
    \item \emph{Interpretation of !!{predicate}s:} the interpretation
    of a $n$-place !!{predicate}~$R$ associates each world $w$ with a
    relation over \emph{the domain of $w$}: $\mAssign{R}{M}[w]\in
    \mDomain{M}(w)$. 
  \item \emph{Interpretation of !!{function}s:} the interpretation 
    an $n$-place !!{function}~$f$ associates each world $w$ with a 
    function \emph{that maps any sequence of objects
    containing some object not in $w$ to an object outside the domain 
    of $w$}. That is, $\mAssign{f}{M}(w)(m_1,\ldots,m_n)\in \mDomain{M}(w)$
    only if $m_1,\ldots,m_n\in \mDomain{M}(w)$.
  \end{enumerate}
\end{defn}

\begin{explain}
  The first constraint ensures that, at a given world, !!a{predicate}
  only applies to objects existing at that world. The second ensures
  that, at a given world, the value of a function is an object of 
  that world only for arguments that are objects existing 
  at that world. 
\end{explain}

\begin{rem} % digress?
  For the negative variant, another approach could have 
  been used. Standard models for (non-modal) free logic use 
partial functions: interpretations of !!{constant}s 
that assign an object to some but not necessarily all !!{constant}s, 
interpretation of !!{function}s that give a value for 
some but not necessarily all constants. We could have 
adopted used a similar approach here. There are a 
few ways of implementing it. One is to let defined
!!{constant}s denote the same object at all worlds, even 
those where that object doesn't exist, but restrict
the identity clause to ensure that $\Obj a=\Obj a$ holds 
only at worlds the object (if any) denoted by $\Obj a$ 
exists. Another is to let constant denotes a single 
object, but only relative to worlds where that object
exist. We leave it to the reader to work out the details.
\end{rem}

\begin{rem}[Model types]
  As in modal logic, a model $\mModel{M}$ is called a $\Ax{D}$-model, 
  $\Ax{T}$-model, etc., if its accessibility relation is serial,
  reflexive, etc., as in \olref[qml][sem][fdm]{tab:five-types}. 
\end{rem}

\end{document}